% Part: incompleteness
% Chapter: theories-computability
% Section: introduction

\documentclass[../../../include/open-logic-section]{subfiles}

\begin{document}

\olfileid{inc}{tcp}{int}

\olsection{تعارف}

\begin{editorial}
ایس حصے نوں فیر توں لکھن دی لوڑ اے۔
\end{editorial}

ساڈے کول ایہ تعریفاں تے قضیے نیں۔
\begin{enumerate}
\item ایس گل دی تعریف کہ کسے فنکشن دا $\Th{Q}$ وچ قابلِ نمائندگی
  ہونا کیہہ اے (\olref[req][int]{defn:representable-fn})
\item ایس گل دی تعریف کہ کسے تعلق دا $\Th{Q}$ وچ قابلِ نمائندگی
  ہونا کیہہ اے (\olref[req][rel]{defn:representing-relations})
\item اک قضیہ جیہڑا آکھدا اے کہ $\Th{Q}$ وچ قابلِ نمائندگی فنکشن
  عین اوہی نیں جیہڑے قابلِ حساب نیں (\olref[req][int]{thm:representable-iff-comp})
\item اک قضیہ جیہڑا آکھدا اے کہ $\Th{Q}$ وچ قابلِ نمائندگی تعلق
  عین اوہی نیں جیہڑے قابلِ حساب نیں
  % PNB-CMP-066: اصل چوتھے حوالے دی غائب کھلّن والی قوس درست۔
  (\olref[req][rel]{thm:representing-rels})
\end{enumerate}

{\em نظریہ} جملیاں دا اوہ سیٹ اے جیہڑا اخذ دے لحاظ نال بند ہووے۔
یعنی جدوں وی $T$، $!A$ نوں ثابت کردا اے، تاں $!A$، $T$ وچ ہوندا اے۔
نظریے نوں سمجھّن دا سوکھا طریقہ ایہ اے کہ جملیاں دا اک مجموعہ
تے اوہ ساریاں گلاں وی منّ لئیے جیہڑیاں انہاں جملیاں توں لازم آندیاں نیں۔
اگے اسیں $\Th{Q}$ نال اوہ {\em نظریہ} مراد لواں گے جیہڑا
\olref[req][int]{sec} دے اٹھ مسلّمیاں توں اخذ ہو سکّن والے سارے
جملیاں دا سیٹ اے۔ یاد رکھو کہ اسیں $\Th{Q}$ دے فارمولیاں نوں عدداں
نال کوڈ کر سکدے آں۔ جے $!A$ ایہو جیہا فارمولا اے، تاں $\Gn{!A}$
نال اوہ عدد مراد لئیے جیہڑا $!A$ نوں کوڈ کردا اے۔ ایس کوڈ بندی
دے حساب نال ہن اسیں پچھ سکدے آں کہ فارمولیاں دے وکھ وکھ سیٹ
قابلِ حساب نیں یا نہیں۔

\end{document}

